\subsection{Acoplamento base--manipulador}
\label{sec:acoplamento_base_manipulador}

As variáveis são para o acoplamento dos movimentos entre a base e o MR:
\begin{equation}
\label{eq:nu_def}
\nu
\;=\;
\begin{bmatrix}
{}^{0}\!\vec v_{B} \\[2pt]
{}^{0}\!\vec \omega_{B} \\[2pt]
\dot{\vec q}
\end{bmatrix},
\end{equation}
onde ${}^{0}\!\vec v_{B}$ e
${}^{0}\!\vec \omega_{B}$
são as velocidades linear e angular da base expressas em $\{0\}$, e
$\dot{\vec q}$ são as velocidades das juntas.

\subsection{Jacobiano composto da base e do MR}
A velocidade cartesiana da ferramenta no mesmo sistema de coordenadas $\{0\}$ passa a ter duas contribuições, sendo elas:

\begin{equation}
\label{eq:xdot_composto}
{}^{0}\!\dot{\vec x}_{E}
\;=\;
\underbrace{\mathbf{J}_{\!B}(\cdot)}_{\text{efeito da base}}\,
\begin{bmatrix}
{}^{0}\!\vec v_{B} \\
{}^{0}\!\vec \omega_{B}
\end{bmatrix}
\;+\;
\underbrace{\mathbf{J}(\vec q)}_{\text{efeito das juntas}}\,\dot{\vec q},
\end{equation}

onde $\mathbf{J}_{\!B}(\cdot)$ depende apenas da posição relativa entre a base $\{B\}$ e
o efetuador $\{E\}$, enquanto que
$\mathbf{J}(\vec q)$ é o Jacobiano geométrico do manipulador.

\subsection{Matriz de inércia e o acoplamento}
Considerando como base livre, as equações de movimento podem ser descritas por:

\begin{equation}
\label{eq:EOM_particionada_compacta}
\mathbf M_{\text{tot}}(\vec q\, )\,\dot{\nu}
\;+\;
\mathbf C_{\text{tot}}(\vec q,\nu)\,\nu
\;=\;
\vec\tau_{\text{tot}}
\;+\;
\vec w_{\text{ext}}.
\end{equation}

Detalhando cada termo, obtém-se:
\begin{equation}
\label{eq:Mtot_def}
\mathbf M_{\text{tot}}(\vec q)
=
\begin{bmatrix}
\mathbf M_{BB}(\vec q) & \mathbf M_{BM}(\vec q) \\
\mathbf M_{MB}(\vec q) & \mathbf M_{MM}(\vec q)
\end{bmatrix}
\end{equation}

\begin{equation}
\label{eq:Ctot_def}
\mathbf C_{\text{tot}}(\vec q,\nu)
=
\begin{bmatrix}
\mathbf C_{BB}(\cdot) & \mathbf C_{BM}(\cdot) \\
\mathbf C_{MB}(\cdot) & \mathbf C_{MM}(\cdot)
\end{bmatrix}
\end{equation}

% Vetores empilhados (opcional, ajuda a reduzir largura na EOM)
\begin{equation}
\label{eq:stacked_defs1}
\dot{\nu} =
\begin{bmatrix}
{}^{0} \dot{\vec v}_{B} \\
{}^{0} \dot{\vec \omega}_{B} \\
\ddot{\vec q}
\end{bmatrix},
\qquad
\nu =
\begin{bmatrix}
{}^{0} \vec v_{B} \\
{}^{0} \vec \omega_{B} \\
\dot{\vec q}
\end{bmatrix},
\end{equation}

% Vetores empilhados (opcional, ajuda a reduzir largura na EOM)
\begin{equation}
\label{eq:stacked_tau}
\vec\tau_{\text{tot}} =
\begin{bmatrix}
\vec \tau_{B} \\
\vec \tau
\end{bmatrix},
\end{equation}

\begin{equation}
\label{eq:stacked_wrench}
\vec w_{\text{ext}} =
\begin{bmatrix}
{}^{\text{ext}}\vec w_{B} \\
{}^{\text{ext}}\vec w_{E}
\end{bmatrix}.
\end{equation}

A equação \eqref{eq:EOM_particionada_compacta} descreve a dinâmica acoplada base--manipulador.
A matriz de inércia \eqref{eq:Mtot_def} agrega a inércia translacional e rotacional da base, a inércia no espaço das juntas e os termos de acoplamento dinâmico entre base e braço; além disso, é simétrica e definida positiva.

A matriz \eqref{eq:Ctot_def} reúne os termos de Coriolis/centrífugos do sistema base e manipulador, incluindo efeitos cruzados.
% Deve satisfazer a propriedade estrutural $\dot{\mathbf M}_{\text{tot}}-2\,\mathbf C_{\text{tot}}$ anti-simétrica, assegurando coerência energética.

Os vetores $\nu$ e $\dot{\nu}$ definidos em \eqref{eq:stacked_defs1} são, respectivamente, as velocidades linear e angular da base, além das velocidades articulares; e as correspondentes acelerações.

Os vetores em \eqref{eq:stacked_tau} são as entradas de atuação (esforços na base e torques/força nas juntas).
Por sua vez, \eqref{eq:stacked_wrench} explicita as excitações externas (força--momento na base e na ferramenta), todas expressas no mesmo referencial e referidas aos pontos $B$ e $E$, respectivamente.

\subsubsection{Momento angular}
Para descrever a interação do momento angular do sistema em torno do centro de massa, utiliza-se:

\begin{equation}
\label{eq:centroidal}
\vec h_{c}
\;=\;
\mathbf A_{c}(\vec q)\,\nu,
\qquad
\dot{\vec h}_{c} \;=\; \vec w^{\text{ext}}_{c}.
\end{equation}

Na ausência de momentos externos
(${}^{\text{ext}}\vec w_{c}=\vec 0$),
o momento angular no centro de massa é conservado: movimentos das juntas alteram $\vec h_{c}$ e, por reação, geram velocidade/atitude na base.
A matriz $\mathbf A_{c}(\vec q)$ é obtida por composição rígida (corpos articulados) e depende apenas da configuração.